\textbf{Conciliación vigente INC-39.} La supervisión independiente añade cuatro EOLR-1 en extremos ordinarios, cuatro monitores y cuatro EOL a la selección desarrollada aquí: 56 monitores y 128 direcciones, cuatro SLC 40/40/24/24 con 88/88/52/52 mA. No cambia seis VEP, cuatro FRPS, ocho baterías 33 Ah ni cuatro 75 Ah. El controlador pasa a 1,232686/5,776 A y 36,078957 Ah de cribado; cada FRPS con relé a 0,625/0,751 A a 24 V y 18,0751 Ah. Los balances de INC-36 que siguen son la base antes de esos cuatro ramales, sin segunda lista activa. La selección y sus límites de subtensión, autonomía y supervisión están en INC/EXT-39.

\section{Lazos, avisadores y fuentes definidos (INC-36)}
\label{sec:sd4-inc36}
Se sustituye la detección ordinaria de cuatro VLF-500-00-UL por cuatro VEP-A00-1P-UL, con dos repuestos fríos ordinarios para cubrir dos sustituciones simultáneas, además del stock dedicado de bancos; se conservan los dos VEP y el repuesto de bancos de INC-30. La ficha 36106\_07 publica máximos de 0,57/0,59 A a 24 V, sin extrapolarlos a 18 V, evitando usar 410 mA nominales VLF como máximo. Los seis sectores conservan redes y escape independientes; las cuatro redes ordinarias se recalculan en ASPIRE para el modelo elegido, sin validar los 72 m/40 puntos por comparación de longitud máxima. La sustitución no completa la confirmación independiente ni la clasificación de bancos \parencite[p.~2]{lote30-vep36106}.

Se seleccionan cuatro FRPS 10-2829-1-0-01-0-1-01 activos y un repuesto frío, uno por VEP ordinario F01/F02/F03/F06. Cada FRPS aporta una salida auxiliar continua no rearmable y mantiene falla individual al panel; 0,57 A externos en reposo son inferiores a 1 A total permitido. Con consumo propio 0,035/0,141 A, cada fuente demanda 0,605/0,731 A DC y el cribado $1,2[24(0,605)+(5/60)(0,731)]=17,4971$ Ah. Se disponen ocho baterías 02-3468 de 12 V/33 Ah activas, dos en serie por fuente, cuatro gabinetes 10-2154-R y reservas previas de batería/gabinete. El conjunto de cuatro fuentes suma 58,08/70,176 W DC; no es potencia AC. Sustituye las dos fuentes y cuatro baterías activas de INC-33. Entrada 240 V/50 Hz, cargadores, pérdidas, protección y autonomía útil se reconcilian por fuente; no se suman Ah en serie ni se atribuye eficiencia no publicada \parencite[pp.~1--2]{lote33-frps}.

Se individualizan 52 monitores 55-046: dieciocho señales de aspirantes, cuatro fallas FRPS, seis abortos, seis presiones, seis IVOS y doce finales de carrera de compuertas. Se seleccionan 54 relés 55-048 (seis de actuación y 48 contactos de monitor auxiliar), seis 55-053 y seis 20-1343. Los cuatro sensores/base 63-1059/63-1060 ordinarios se conservan; los dos confirmadores de bancos siguen sin variante apta seleccionada. La reserva pasa justificadamente de 122 a 124 direcciones por dos fallas FRPS adicionales. Una entrada identifica un estado, sin agrupar presión e IVOS; relés hacia lectores auxiliares son contactos secos, sin alimentar sus cargas desde SLC.

El límite SLC es 100 mA por lazo, además de 254 direcciones. La planilla da 2 mA de alarma por monitor, relé, pulsador o sensor, y 6 mA por RCM: la reserva desarrollada alcanza 272 mA. Dos lazos no bastan aunque quepan las direcciones. Se selecciona un SLM 10-2473 activo y uno de reserva: SLC1 F02/F04 y SLC2 F03/F05, cuarenta direcciones/88 mA cada uno; SLC3 F01 y fallas FRPS F01/F02, y SLC4 F06 y fallas FRPS F03/F06, veintidós/48 mA cada uno. Cada sector reserva veinte direcciones: sensor, tres señales aspirante, RCM, manual, aborto, presión, IVOS, dos compuertas, relé de actuación y ocho relés auxiliares. Los dos sensores de bancos se presupuestan sólo provisionalmente con 2 mA; su elección puede cambiar el balance. Se conserva separación física lateral A/B y retorno Class A, con aislamiento por dispositivo compatible; la estación 20-1343 no tiene versión aisladora y no se declara Class X de todo el circuito. Resistencia/capacitancia y protección de ambos recorridos se verifican con trazado real \parencite[pp.~3-5, 3-8 y ap.~A]{lote33-cheetah11}.

Se seleccionan seis System Sensor P2RL-SP activos y uno de reserva, avisador de pared rojo de dos hilos con inscripción FUEGO, compatible según Fike 06-186 p.~26. Un equipo por acceso de sector se instala en espacio ordinario seco; los de bancos van fuera de su recinto clasificado. Se usa presupuesto de 262 mA por equipo en DC filtrada, peor promedio publicado por Honeywell; los 326 mA corresponden a FWR, que no se usa en el presente cálculo. P10 alimenta F01/F02/F04 y P11 F03/F05/F06, tres equipos/0,786 A por circuito, por debajo de 2 A. Ambos NAC se programan con sincronización System Sensor y alarma de evacuación general; el estado y zona de predescarga/descarga se identifican además en panel/RDU2, no mediante dos tonos supuestos del mismo equipo. La cobertura acústica/visual y selección de candela/tono se reciben mediante cálculo y medición; la cantidad por acceso no acredita audibilidad en todo el recinto \parencite[secs.~3.1--3.3 y p.~26]{lote36-fike186}\parencite[p.~11, tabla~9]{lote36-honeywell54504}.

NAC P10/P11 se cablean Class B con un EOL 10-2570 de 1,2 k$\Omega$ en el último equipo, dos activos y reserva por tipo; no se coloca EOL junto al panel. Las 52 entradas Class B usan EOL 10-2625 de 39 k$\Omega$ en extremo supervisado; resistor serie 10-2530 de 14 k$\Omega$ donde el esquema requiera distinguir cortocircuito, sin sustituir un valor por otro. Se preservan supervisión de falla eléctrica, inversión de polaridad NAC y disparo independiente. El balance SLC incorpora corrientes de planilla con terminación prescrita; no se añade $24/39000$ indiscriminadamente a todos los módulos como si fuera una carga DC auxiliar permanente \parencite[secs.~3.9.2, 4.16.6 y ap.~C]{lote33-cheetah11}.

\section{Balance por fuente y caída de actuación (EXT-36)}
\label{sec:sd4-ext36}
Se asignan al controlador Cheetah un VEP de banco, tres RCM/DFIA, SLM, ambos lazos principales y de expansión, los 124 dispositivos presupuestados, RDU2 y seis avisadores. Al SPS se asignan el segundo VEP y tres RCM/DFIA. El cálculo conservador usa RCM auxiliar de 0,0364/0,010 A además de DFIA 0,0254/0,930 A, pendiente cotejo OEM de supervisión para evitar doble conteo. SLC suma 0,064346/0,272 A; controlador 0,275 A, SLM 0,100 A, SPS 0,040 A, RDU2 0,036/0,139 A. Resultados: controlador 1,230746/5,768 A y SPS 0,795400/3,450 A, inferiores respectivamente a 1,5/6 A. Total 2,026146/9,218 A, inferior a 3/12 A. No se cuenta un repuesto frío como carga activa. Las cargas de bobinas finales de interlock se verifican en sus fuentes propias y no se esconden detrás del consumo del contacto de relé \parencite[ap.~A]{lote33-cheetah11}.

Las dos salidas continuas del controlador se asignan: una a dos DFIA/RCM (1,88 A de alarma); otra a un DFIA/RCM, VEP y RDU2 (1,669 A). Las tres del SPS: dos DFIA/RCM (1,88 A), uno (0,94 A) y VEP (0,59 A), respectivamente. Cada salida permanece bajo 2 A. Los NAC conservan bornes propios; los actuadores reciben 24 V externos al SLC y uno por RCM. La liberación no se alimenta del contacto seco ni de una FRPS de aspiración. El presupuesto adopta alarma simultánea de todas las cargas como extremo eléctrico aunque una sola alarma no ordene seis descargas.

Con 24 h de reposo, cinco minutos de alarma y factor 1,2, la capacidad propia calculada es 36,0223 Ah para controlador y 23,2525 Ah para SPS. Se conservan los dos pares 75 Ah 02-4206 y dos gabinetes 10-2236-R, uno por conexión OEM, sin paralelo de obra. Se comprueban temperatura, curva útil, tensión de corte, recarga y cable de batería; no basta comparar esos Ah con etiqueta 75 Ah. Gabinetes batería/SPS permanecen a distancia permitida por instrucciones OEM, fuera de bancos, sin bloquear acceso ni compartir soporte A/B de potencia. Las longitudes hasta actuadores exteriores requieren reconciliar ubicación de panel, gabinetes y canalizaciones; no se consideran medidas por disponer dimensiones de sala.

Para la salida que alimenta dos DFIA se desarrolla un límite de ruta, no una medición: cobre multifilar recubierto AWG14, resistencia 10,7 $\Omega$/km a 20 grados C según tabla del manual, longitud máxima de proyecto 20 m de ida hasta el último actuador, factor propio 1,13755 a 55 grados C y 0,20 V de reserva para conectores/RCM. A 1,88 A, $\Delta V=2(0,020)(10,7)(1,13755)(1,88)+0,20=1,1153$ V. Una salida mínima de 22,0 V dejaría 20,8847 V, sobre los 20,4 V FM. El manual exige 20,4 V en cualquier condición pero no acredita aquí mínimo de 22 V al final de batería; por ello esta desigualdad no cierra la caída ni la autonomía. Se verifica mínimo real de fuente y ruta antes de aceptar AWG14; una ruta más larga no se autoriza con este presupuesto. Las rutas NAC y aspirantes tienen su comprobación independiente de tensión mínima, máximo consumo y sensibilidad a avería \parencite[pp.~3-4, 4-3 y ap.~A nota~12]{lote33-cheetah11}.

\section{Servicios admitidos del paso y condición gas/fuego (PEN-36)}
\label{sec:sd4-pen36}
Se conserva la solución de treinta aberturas y sesenta conjuntos de ambas caras PEN-33. La brecha puntual es el cotejo de cubierta, conductores y diámetro real: C-AJ-8309 identifica combinaciones PVC, EPR/PVC y XLPE/PVC, con distintos T; la fibra PVC hasta 6 mm no acredita OS2 LSZH ni una designación genérica de baja emisión de humo. No se sustituye la cubierta de enlaces de CER-4 sólo para forzar un listado: se exige cable concreto con ficha, reacción al fuego, compatibilidad de enlace y alcance del paso, y módulos por cada cara. Para control AWG14 se conserva el calibre de cálculo sin inventar cubierta/diámetro comercial. El listado debe incluirlo antes de cerrar compra y empaquetado; no se repite la masa como evidencia de aptitud del cable \parencite[pp.~3--4]{lote30-caj8309}.

La elección eléctrica no elimina las brechas de incendio de bancos. Confirmador independiente apto, clasificación normal/falla/posaislamiento, certificado íntegro VEP y accesorios, ASPIRE, emisión residual, retención/alivio, red hidráulica y concentración por riesgo conservan sus verificaciones. Se mantienen retirada de recarga ante pérdida de caudal/señal válida y la purga INT-31 como cribado sin temporizador habilitante. No se declara extinción automática de bancos satisfecha mediante inhibición permanente o simple corte AC.

El inventario integrado fija seis reservas compatibles de 55-046 y seis de 55-048, aplicando $\lceil0,10N\rceil$ a 52 y 54 activos; seis EOL 39 k$\Omega$ de reserva para 52 activos. Sustituyen el stock previo de dos monitores. No se cuentan como carga energizada.

La salida auxiliar FRPS declara 18,6--28,5 V, 1 A y ausencia de supervisión propia. El contacto de falla de la fuente no acredita supervisión de toda la rama del detector. Frente al mínimo VEP de 18 V, el extremo 18,6 V deja sólo 0,6 V de caída admisible; consumo en ese extremo, tendido y supervisión del ramal requieren verificación específica. Los 0,57/0,59 A documentados a 24 V no son máximos acreditados a 18,6 V. Esa condición conserva pendiente alimentación/autonomía de aspiración.\parencites[pp.~1--2]{lote33-frps}[p.~2]{lote30-vep36106}.
